% §7 Discussion and Conclusion
%
% Restructured to center on the star-tree theorem's implications.

%%====================================================================
\subsection{Summary}

We have proved that the Seeley--Richard--Love index-set construction
satisfies the canonical anticommutation relations if and only if the
underlying tree is a star (\cref{thm:star-tree}).  The proof
identifies a generic structural failure — a parity-set cancellation
along any depth-$\ge 2$ path — that is invisible when encodings are
studied individually.  This result classifies the space of $n^{n-1}$
encoding trees into three construction-compatibility regions and
establishes the path-based construction (Construction~B) as the only
universal method.

%%====================================================================
\subsection{Implications for the literature}

\paragraph{The SRL framework is narrower than cited.}
The index-set framework of Seeley, Richard, and
Love~\cite{seeley2012} is widely cited as a unified treatment of
fermion-to-qubit encodings.  Our result shows that this unification
covers only star trees --- depth-1 graphs that produce relabelled
Jordan--Wigner or Parity encodings.  The Bravyi--Kitaev encoding,
commonly presented as an instance of the SRL framework, actually uses
a separate Fenwick-specific construction (Construction~F) that is
algebraically distinct from the generic algorithm.  We recommend that
future presentations distinguish these constructions explicitly.

\paragraph{Silent encoding errors.}
The failure mode identified in \cref{sec:st-failure} is particularly
concerning because it is silent: Construction~A applied to a non-star
tree produces well-formed Pauli strings with correct widths and
plausible coefficients.  The encoded operators pass superficial checks
(correct number of terms, correct register size) but fail the CAR,
producing Hamiltonians with wrong eigenspectra.  The diagnostic
$D(T)$~\eqref{eq:diagnostic} provides a simple test that
implementations should apply when deriving index sets from novel tree
structures.

\paragraph{Retroactive justification for path-based constructions.}
The work of Jiang et al.~\cite{jiang2020} and Miller et
al.~\cite{miller2023bonsai} introduced path-based tree encodings as
a generalization of Bravyi--Kitaev.  Our result provides a
retroactive explanation of why this construction was necessary: the
alternative (deriving index sets from the tree) fails for all but
the simplest trees.  The universality of Construction~B is not just
convenient but \emph{required}.

%%====================================================================
\subsection{Limitations}

\paragraph{Computational proof scope.}
The exhaustive enumeration in \cref{sec:val-census} covers
$n \le 6$.  The structural proof (\cref{thm:star-tree}) provides a
general argument for all~$n$, but relies on identifying a specific
failure witness (the $\{d_w, c_v\}$ anticommutator for a depth-2
path).  We have not proved that this is the \emph{only} failure mode,
though the enumeration data are consistent with it being generic.

\paragraph{No computational speedup.}
The encoding pipeline has $O(n^5)$ complexity, matching existing
implementations~\cite{mcclean2020, bergholm2022}.  The advantage is
structural (correctness by construction, exact phases, extensibility
to mixed systems) rather than asymptotic.

\paragraph{Bosonic exactness.}
Algebraic exactness (\cref{thm:exactness}) applies only to the
fermionic pipeline.  Bosonic encodings introduce $\sqrt{n+1}$
coefficients during truncated-space decomposition.

%%====================================================================
\subsection{Open questions}

\begin{enumerate}
  \item \textbf{Bijective proof of $|M(n)| = (n{-}1)!$.}
    The identity between monotonic tree counts and shifted factorials
    is verified computationally (\cref{prop:monotonic-count}).  A
    constructive bijection via heap-labeled tree
    correspondences~\cite{bergeron1992} would be a useful
    contribution to enumerative combinatorics in the encoding context.

  \item \textbf{Problem-adapted tree optimization.}
    Given a molecular Hamiltonian, what tree topology minimizes the
    total Pauli weight of the encoded Hamiltonian?  The framework
    enables systematic exploration, but the optimization landscape
    is combinatorially vast ($n^{n-1}$ candidates).

  \item \textbf{Higher-arity tree encodings.}
    The branching factor~$\le 3$ restriction arises from the three
    non-identity single-qubit Paulis.  Multi-qubit link labels could
    support higher arity, at the cost of the one-qubit-per-mode
    property.

  \item \textbf{Bosonic tree encodings.}
    Can the path-based construction be adapted to truncated bosonic
    modes, yielding a unified tree-based encoding for mixed systems?
\end{enumerate}

%%====================================================================
\subsection{Conclusion}

The star-tree theorem corrects a widespread assumption about the
generality of index-set fermion-to-qubit encodings.  By making the
parameterization explicit, we identified a structural limitation that
is invisible when encodings are studied one at a time.  The three
constructions (A, B, F) and their distinct domains provide a cleaner
conceptual framework for the encoding landscape.  The path-based
construction emerges as the only universal method, and its provably
optimal instance (balanced ternary tree, $O(\log_3 n)$) resides
outside the domain of the index-set approach.
